% SOURCE_AUTHORITY: sga5_fr_workpass.tex, lines 432--451 (LF-normalized unit).
\subsection*{Observações 1.8}
\begin{enumerate}[label=\alph*)]
\item Mediante uma redução padrão (truncando $F$ e usando que $\D_c(X)$ é triangulada), vê-se que a condição (ii) equivale a

(ii~bis): para todo $A_X$-módulo construtível $F$ e todo $i \in \Z$, $\underline{\Ext}^i(F,K)$ é construtível.

\item A condição (ii~bis) implica que $K \in \ob \D_c^{+}(X)$. Ao contrário do que ocorre com os feixes coerentes (EGA~$0_{\mathrm{III}}$~12.3.3), talvez a recíproca não se verifique automaticamente. Todavia, ela se verifica (ver n.\textsuperscript{o}~3) nos casos ``bons'', por exemplo quando se dispõe, sobre $X$, da resolução de singularidades e da pureza (em particular, se $X$ é um esquema excelente de característica zero).

\item Mediante uma redução sobre $F$, vê-se que a condição (iii) equivale a (iii~bis): todo $A_X$-módulo construtível $F$ é reflexivo com respeito a $K$.

\item Se $K \in \ob \D^{+}(X)$ satisfaz (i) e (ii) (respectivamente, é dualizante), o funtor $\underline{D}_K$ induz um funtor (respectivamente, uma equivalência de categorias) $(\D_c^{b}(X))^{\circ} \to \D_c^{b}(X)$.

\item Se $K \in \ob \D^{+}(X)$ tem dimensão \emph{injetiva} finita (cf.~1.5.2) e satisfaz (ii), então, para todo $F \in \ob \D_c(X)$, tem-se $\underline{D}_K(F) \in \D_c(X)$, e $\underline{D}_K$ induz um funtor $(\D_c^{+}(X))^{\circ} \to \D_c^{-}(X)$.

Se, além disso, $K$ é dualizante, todo $F \in \ob \D_c(X)$ é reflexivo em relação a $K$, e $\underline{D}_K$ induz equivalências de categorias
\[
(\D_c(X))^{\circ} \xrightarrow{\approx} \D_c(X), \quad (\D_c^{-}(X))^{\circ} \xrightarrow{\approx} \D_c^{+}(X), \quad (\D_c^{+}(X))^{\circ} \xrightarrow{\approx} \D_c^{-}(X).
\]
Ignora-se se esta conclusão continua válida sem a restrição de que $K$ tenha dimensão injetiva finita.
\end{enumerate}
